% Source pp. 93-101 (Lectures7-12.pdf).
\section{Diagram Evaluation. Divergences}

Now we turn to actual evaluation of the diagram contributions. For this purpose
in most cases it is much more convenient to use the Feynman rules in the
momentum space, evaluating the Fourier transformed correlation functions
\begin{multline}\label{eq:FTW}
\int W^{(n)}(x_1,\dots,x_n)\,e^{-ip_1\cdot x_1}\cdots e^{-ip_n\cdot x_n}\dd^4x_1\cdots \dd^4x_n\\
=(2\pi)^4\,\delta^{(4)}(p_1+\dots+p_n)\,\tilde W^{(n)}(p_1,\dots,p_n)\,.
\end{multline}
Here again $W^{(n)}$ denotes \emph{connected} $n$-point correlation function.
The delta function in the r.h.s.\ appears because of the translation
invariance: $W^{(n)}(x_1,\dots,x_n)$ depends only on the differences
$x_i-x_j$.

I will skip detailed derivation of the diagram rules in the momentum space
(look up \textbf{Sect.~4.4 of PS} for the details). Very briefly, in any
diagram one uses the Fourier transform
\[
D(x-x')=\int\frac{\dd^4p}{(2\pi)^4}\,\tilde D(p)\,e^{ip\cdot(x-x')}\,,\qquad
\tilde D(p)=\frac1{p^2+m^2}
\]
of all propagators. Then the integral over all $x$ associated with vertices, as
well as over all ``external'' points $x_i$, can be evaluated explicitly yielding
the momentum delta-functions. As the result one obtains the

\subsection*{Momentum-space Feynman rules:}
\begin{enumerate}[label=(\roman*)]
\item ``Internal'' propagators carry (directed) momentum to be integrated over,
\[
\dia{\draw[prop] (-1.5,0) -- (1.5,0); \draw[arr] (-0.5,0.35) -- (0.5,0.35) node[above,pos=0.5] {\small$p$};}
=\dia{\draw[prop] (-1.5,0) -- (1.5,0); \draw[arr] (0.5,0.35) -- (-0.5,0.35) node[above,pos=0.5] {\small$-p$};}
=\int\frac{\dd^4p}{(2\pi)^4}\,\frac1{p^2+m^2}\,.
\]
\item ``External'' legs carry fixed momenta
\[
\tilde W^{(n)}(p_1,\dots,p_n)=
\diac{\foreach \a/\l in {135/p_1,90/p_2,45/p_3,0/p_4} {\draw[prop] (0,0) -- (\a:2.2);
  \draw[arr] (\a:2.0) ++(\a+90:0.35) -- ++(\a+180:0.8); \node at ($(\a:2.2)+(\a-90:0.5)$) {\small$\l$};}
  \node[circle,draw,fill=white,inner sep=1.3pt] at (0,0) {};
  \draw[densely dotted] (165:1.4) arc[start angle=165,end angle=345,radius=1.4];}\,,
\]
\item The vertices conserve the momentum
\[
\diac{\foreach \a/\l in {135/p_1,45/p_2,315/p_3,225/p_4} {\draw[prop] (0,0) -- (\a:2);
  \draw[arr] (\a:1.9) ++(\a+90:0.3) -- ++(\a+180:0.7); \node at ($(\a:1.9)+(\a+90:0.75)$) {\small$\l$};}
  \node[vtx] at (0,0) {};}
=-\lambda\,(2\pi)^4\,\delta^{(4)}(p_1+\dots+p_4)\,,
\]
\item Integration $\int\frac{\dd^4p_i}{(2\pi)^4}$ is performed over all the
momenta associated with the ``internal'' propagators.
\end{enumerate}

\subsection*{Some properties of \texorpdfstring{$\tilde W^{(n)}$}{W(n)}}

In the momentum representation many operations simplify. For example the
``amputation'' takes very simple form
\[
\tilde W^{(n)}(p_1,\dots,p_n)=\left[\prod_{i=1}^n\tilde W(p_i)\right]\,\tilde W^{(n)}_{\text{amp}}(p_1,\dots,p_n)\,,
\]
where $\tilde W(p)$ is the Fourier transform of $W(x-x')$. Similar
simplification occurs to the relation to proper vertices.

\begin{exercise}
Express $\tilde W^{(6)}(p_1,\dots,p_6)$ in terms of
$\tilde\Gamma^{(6)}(p_1,\dots,p_6)$, $\tilde\Gamma^{(4)}$ and $\tilde W$.
\end{exercise}

The momentum space $\tilde\Gamma^{(2)}(p)$ is just the inverse of
$\tilde W(p)$ in usual sense,
\[
\tilde\Gamma^{(2)}(p)=\frac1{\tilde W(p)}=p^2+m^2+\tilde\Sigma(p)\,,
\]
where $-\tilde\Sigma(p)$ is the sum of all one-particle irreducible two-leg
diagrams in the momentum space.

\subsection*{Leading contribution to \texorpdfstring{$\tilde\Sigma$}{Sigma}. Mass renormalization}

Let us calculate some perturbative contributions to $\tilde\Sigma(p)$. In the
leading order we have
\[
\tilde\Sigma(p)=-\Bigl(\,\dia{\draw[prop,densely dashed] (-2,0) -- (2,0); \draw[prop] (0,0.6) circle (0.6);
  \draw[arr] (-0.3,1.45) -- (0.3,1.45); \node[above] at (0,1.45) {\small$k$}; \node[vtx] at (0,0) {};}\,\Bigr)
=\frac\lambda2\int\frac{\dd^4k}{(2\pi)^4}\,\frac1{k^2+m^2}
=\frac{\lambda}{16\pi^2}\int_0^\infty\frac{|k|^3\dd|k|}{|k|^2+m^2}\,.
\]
This integral is badly divergent as $k\to\infty$. What should we do about it?
We can get rid of the meaningless expression by introducing a ``cutoff'', say,
replacing
\[
\frac1{k^2+m^2}\ \to\ \frac1{k^2+m^2}\,\Phi(k^2/\Lambda^2)\,,
\]
with the ``cutoff factor'' $\Phi(\xi)$ decaying at $\xi\to\infty$ fast enough
to make the integral convergent. Then
\[
\frac\lambda2\int\frac{\dd^4k}{(2\pi)^4}\,\frac1{k^2+m^2}\,\Phi(k^2/\Lambda^2)=\frac\lambda2\,m^2\,F(\Lambda^2/m^2)\,.
\]
Just as in the case of KG vacuum energy, this move alone brings no real
progress: it replaces something meaningless by something undetermined.

In this simple diagram one can note that the contribution has no momentum
dependence, and enters the quantity $\tilde\Gamma^{(2)}(p)$ in the combination
\[
\tilde\Gamma^{(2)}(p)=p^2+m^2+\frac\lambda2\,m^2\,F(\Lambda^2/m^2)+O(\lambda^2)\,.
\]
We see that no matter what $\Phi$ and $\Lambda$ are, the parameter $m$ in our
action is not actual mass of the particle; the interaction leads to a shift
that happens to depend on the cutoff. (We will explain later that actual mass
$m_{\mathrm{ph}}$ of the particle in interacting theory is determined through
the equation $\tilde\Gamma^{(2)}(p)\big|_{-p^2=m^2_{\mathrm{ph}}}=0$.)
Therefore it is reasonable to change the notations, denoting $m_0^2$ the
coefficient in front of $\varphi^2$ in the action. Then, up to terms
$\sim\lambda^2$ and higher
\[
\tilde\Gamma^{(2)}(p)=p^2+\underbrace{m_0^2+\frac\lambda2\,m_0^2\,F(\Lambda^2/m_0^2)}_{m^2}
\]
Recall that we have already observed similar phenomenon in our path integral
representation of the free propagator $D$. We started with a discrete
approximation of the path integral (an analog of the cutoff), and we were
forced to take some parameter $m_0$ instead of $m$ in the cutoff action
\[
\mathcal A=m_0\,\text{Length}\,;
\]
we then found that $m_0$ had to be given some dependence on the cutoff
parameter $\Delta$ (which plays the same role as $\Lambda^{-1}$ here) to keep
actual mass $m$ finite in the limit $\Delta\to0$.

Similar idea can be applied here. Assume that the parameter $m_0^2$ in the
action ($m_0$ is usually called the ``bare mass'') depends on $\Lambda$ in such
a way that the actual mass
\[
m^2=m_0^2+\frac\lambda2\,m_0^2\,F(\Lambda^2/m_0^2)+O(\lambda^2)
\]
remains finite as $\Lambda\to\infty$. Then, to this order,
\[
\tilde\Gamma^{(2)}(p)=p^2+m^2+O(\lambda^2)
\]
has finite limit. The result is expressed in terms of the actual mass; all
ambiguities associated with the cutoff are ``absorbed'' by the ``bare'' mass
parameter.

This idea can be refined as follows. We start with the action
\[
\mathcal A=\int\left[\frac12\,(\partial\varphi)^2+\frac{m_0^2}2\,\varphi^2+\frac\lambda{4!}\,\varphi^4\right]\,.
\]
Let us write
\[
m_0^2=m^2+\delta m^2\,,
\]
where $m^2$ is assumed to coincide with actual mass squared; at this point the
term $\delta m^2$ is just our way to acknowledge that the mass parameter in the
action need not be equal to $m^2$. We then split the action into the free and
interaction parts as follows
\[
\mathcal A=\int\left[\underbrace{\frac12\,(\partial\varphi)^2+\frac{m^2}2\,\varphi^2}_{\mathcal A_0}
+\underbrace{\frac{\delta m^2}2\,\varphi^2+\frac\lambda{4!}\,\varphi^4}_{\mathcal A_I}\right]\,,
\]
Now the propagator is
\[
\dia{\draw[prop] (-1.5,0) -- (1.5,0); \draw[arr] (-0.5,0.35) -- (0.5,0.35) node[above,pos=0.5] {\small$p$};}
=\frac1{p^2+m^2}\,,
\]
not $\frac1{p^2+m_0^2}$, as it was in the original perturbation theory, but we
have an additional two-leg vertex
\[
\dia{\draw[prop] (-1.8,0) -- (1.8,0); \node[xvtx] at (0,0) {};
  \draw[arr] (-1.4,0.35) -- (-0.6,0.35) node[above,pos=0.5] {\small$p_1$}; \draw[arr] (1.4,0.35) -- (0.6,0.35) node[above,pos=0.5] {\small$p_2$};}
=-\delta m^2\,(2\pi)^4\,\delta^{(4)}(p_1+p_2)\,,
\]
which we will rather write in a shorter form
\[
\dia{\draw[prop] (-1.8,0) -- (1.8,0); \node[xvtx] at (0,0) {};
  \draw[arr] (-1.4,0.35) -- (-0.6,0.35) node[above,pos=0.5] {\small$p$}; \draw[arr] (0.6,0.35) -- (1.4,0.35) node[above,pos=0.5] {\small$p$};}
=-\delta m^2\,,
\]
reflecting the fact that in this case the momentum delta function is easily
integrated out.

The term $\frac{\delta m^2}2\varphi^2$ in the action is the simplest
\emph{counterterm}, and the associated two-leg vertex is called the counterterm
vertex. In this modified perturbation theory we have (up to terms
$\sim\lambda^2$ and higher)
\begin{align*}
\tilde\Sigma(p)&=-\Bigl(\,\dia{\draw[prop,densely dashed] (-2,0) -- (2,0); \draw[prop] (0,0.6) circle (0.6); \node[vtx] at (0,0) {};}
+\dia{\draw[prop,densely dashed] (-2,0) -- (2,0); \node[xvtx] at (0,0) {};}\,\Bigr)\\
&=\frac\lambda2\int\frac{\dd^4k}{(2\pi)^4}\,\frac1{k^2+m^2}\,\Phi(k^2/\Lambda^2)+\delta m^2=\frac\lambda2\,m^2\,F(\Lambda^2/m^2)+\delta m^2\,.
\end{align*}
Since we insist that $m$ is the actual mass, we have to set
\[
\delta m^2=-\frac\lambda2\,m^2\,F(\Lambda^2/m^2)+O(\lambda^2)\,.
\]
Note that this looks different from the relation
\[
m_0^2=m^2-\frac\lambda2\,m_0^2\,F(\Lambda^2/m_0^2)+O(\lambda^2)
\]
we have obtained within the original perturbation theory; however these two
relations are consistent since the difference is $\sim\lambda^2$.

As we have seen, in the modified perturbation theory, to the order $\sim\lambda$
the counterterm diagram exactly cancels the cutoff dependent bubble diagram,
\[
\dia{\draw[prop,densely dashed] (-2,0) -- (2,0); \node[xvtx] at (0,0) {}; \node[above] at (0,0.2) {\scriptsize$\delta m^2$};}
+\dia{\draw[prop,densely dashed] (-2,0) -- (2,0); \draw[prop] (0,0.6) circle (0.6); \node[vtx] at (0,0) {};}=0\,,
\]
by choice of the parameter $\delta m^2$. The result has no dependence on
$\Lambda$ and $\Phi$. Note also that this cancellation occurs inside more
complicated diagrams, like
\[
\diac{\draw[prop,densely dashed] (-3,0) -- (-1.8,0) (1.8,0) -- (3,0);
  \draw[prop] (-1.8,0) -- (1.8,0); \draw[prop] (1.8,0) arc[start angle=0,end angle=180,radius=1.8];
  \draw[prop] (0,0) arc[start angle=0,end angle=-180,radius=0.9]; \draw[prop] (1.8,0) arc[start angle=0,end angle=-180,radius=0.9];
  \draw[prop] (0,2.4) circle (0.6); \foreach \p in {(-1.8,0),(0,0),(1.8,0),(0,1.8)} \node[vtx] at \p {};}
+\diac{\draw[prop,densely dashed] (-3,0) -- (-1.8,0) (1.8,0) -- (3,0);
  \draw[prop] (-1.8,0) -- (1.8,0); \draw[prop] (1.8,0) arc[start angle=0,end angle=180,radius=1.8];
  \draw[prop] (0,0) arc[start angle=0,end angle=-180,radius=0.9]; \draw[prop] (1.8,0) arc[start angle=0,end angle=-180,radius=0.9];
  \node[xvtx] at (0,1.8) {}; \foreach \p in {(-1.8,0),(0,0),(1.8,0)} \node[vtx] at \p {};}
=0
\]
For this reason the diagrams involving the bubble can be ignored in
calculations, if we keep in mind that their contributions are canceled by the
mass counterterm.

The idea behind this modified perturbation theory can be extended to obtain the
so called \emph{renormalized perturbation theory}, in which all divergences of
original perturbation theory are eliminated. But first let's find out exactly
what diagrams diverge.

\subsection*{Divergences in \texorpdfstring{$\varphi^4$}{phi**4}}

Consider $n$-leg proper vertex
\[
\tilde\Gamma^{(n)}(p_1,\dots,p_n)=
\diac{\foreach \a/\l in {135/p_1,90/p_2,45/p_3,0/p_4} {\draw[prop] (\a:0.8) -- (\a:2.2);
  \draw[arr] (\a:2.0) ++(\a+90:0.35) -- ++(\a+180:0.8); \node at ($(\a:2.2)+(\a-90:0.5)$) {\small$\l$};}
  \draw[prop,pattern=north east lines] (0,0) circle (0.8);
  \foreach \a in {135,90,45,0} \node[vtx,minimum size=4pt] at (\a:0.8) {};
  \draw[densely dotted] (165:1.4) arc[start angle=165,end angle=345,radius=1.4];}\,.
\]
How to characterize diagrams which contribute to this vertex? Assume that the
diagram contains
\[
\begin{array}{lll}
V & \text{vertices} & \dia{\pic{xvertex};}\\[8pt]
P & \text{propagators} & \dia{\draw[prop] (0,0) -- (2,0);}\\[8pt]
I & \text{momentum integrals} & \displaystyle\int\frac{\dd^4k}{(2\pi)^4}\,.
\end{array}
\]
These numbers are not all independent. First, $V$ vertices must absorb exactly
$4V$ lines, $2P$ coming from the propagators, and $n$ from the external
(amputated) legs, i.e.\ $4V=2P+n$, or
\[
2P=4V-n\,.
\]
Next, there is a 4-momentum flowing in each propagator, but there is also a
momentum delta-function associated with each vertex. Each of these
delta-functions reduces the number of 4-momentum integrations by 1. This is
true for all delta-constraints but one, because the momentum-space connected
correlation functions must have one overall delta function left, namely
\[
(2\pi)^4\,\delta^{(4)}(p_1+\dots+p_n)\,.
\]
Therefore
\[
I=P-V+1=\bigl(P=2V-n/2\bigr)=V-n/2+1\,.
\]
Thus we have two equations
\[
\begin{cases}
2P=4V-n\\
4I=4V-2n+4
\end{cases}
\]
Now, the expression for the diagram has the following general form
\[
\int\frac{\dd^4k'_1\cdots \dd^4k'_I}{(q_1^2+m^2)\cdots(q_P^2+m^2)}\,,
\]
where $k'_1,\dots,k'_I$ are the integration momenta left over after eliminating
the delta-functions, and $q_1,\dots,q_P$ are some linear combinations of $k'_j$
and the external momenta $p_i$. For our present analysis it is irrelevant if we
have $m^2$ or $m_0^2$ in the propagators.

We are concerned with possible divergence as $k'_j\to\infty$ (which generally
implies $q_i\to\infty$). Roughly speaking, the above integral diverges if the
overall power of the momenta in the integration measure in the numerator
exceeds or equals the overall power in the denominator. Therefore let us define
the \emph{superficial degree of divergence} $D$ as
\[
D=4I-2P\,.
\]
The integral definitely diverges if $D\ge0$ (but $D<0$ does not guarantee
convergence, see below). If we introduce a cutoff momentum $\Lambda$, the
diagram of the superficial degree of divergence $D$ is $\sim\Lambda^D$.

Substituting the above expressions for $P$ and $I$, we find
\[
D=4V-2n+4-4V+n=4-n\,.
\]
We see that the number of the vertices $V$ cancels, and the superficial degree
of divergence depends on $n$ only. In addition, we see that only finite number
of the proper vertices ($n=0$, $n=2$ and $n=4$) are superficially divergent.

Incidentally, if a perturbation theory has finitely many superficially
divergent vertices (as the $\varphi^4$ perturbation theory does), it is called
\emph{renormalizable}, for the reason which will become clear later. Otherwise,
when infinitely many proper vertices are superficially divergent, we are
dealing with \emph{non-renormalizable} perturbation theory.

\begin{exercise}
Analyze the superficial degree of divergence for all proper vertices in
$\varphi^6$ theory.
\end{exercise}

In the $\varphi^4$ perturbation theory, the proper vertices with $n>4$ are
superficially convergent. This does not mean all the diagrams for these
vertices are given by convergent integrals (that is why we use the name
``superficial'' for that, essentially dimensional, analysis). The reason is
that the diagrams for $\Gamma^{(n)}$ with $n\ge6$ still can have 2- or 4-leg
subdiagrams. For example, the diagram series for $\Gamma^{(6)}$ contains the
diagrams like these
\[
\diac{\coordinate (a) at (-1.5,1.5); \coordinate (b) at (1.5,1.5); \coordinate (c) at (0,-1);
  \draw[prop] (a) -- (b) -- (c) -- (a); \draw[prop] (0,1.5) ellipse (0.55 and 0.32);
  \draw[densely dotted] (-0.9,1.0) rectangle (0.9,2.0);
  \foreach \p/\q in {a/180,a/120,b/0,b/60,c/240,c/300} \draw[prop,densely dashed] (\p) -- ++(\q:0.9);
  \foreach \p in {(a),(b),(c),(-0.55,1.5),(0.55,1.5)} \node[vtx] at \p {};}
\qquad\qquad
\diac{\coordinate (a) at (-1,1); \coordinate (b) at (1,1); \coordinate (c) at (-1,-1); \coordinate (d) at (1,-1);
  \draw[prop] (a) -- (c) -- (d) -- (b); \draw[prop] (0,1) ellipse (1 and 0.3);
  \draw[densely dotted] (-1.5,0.5) rectangle (1.5,1.6);
  \foreach \p/\q in {a/135,b/45,c/180,c/270,d/0,d/270} \draw[prop,densely dashed] (\p) -- ++(\q:0.9);
  \foreach \p in {a,b,c,d} \node[vtx] at (\p) {};}
\]
which contain divergent blocks with two or four legs. We will come back to
decide how to handle such subdiagrams later.
